\section{Definición espectral de la masa como medida de persistencia de \(1\) ruta acoplada a un régimen físico}
La masa es la medida espectral de la persistencia de una ruta cuando una sección genealógica se acopla a un régimen físico. No es un atributo escalar adherido de antemano a un nombre de partícula. La dinámica APP--TRIT--TPK construye la ruta, su memoria, su firma y su carácter; la Mecánica Dimensional realiza esa genealogía en una recta de masa y determina, mediante el acoplamiento, el subespacio observable y su espectro.

Una partícula elemental se representa por una sección persistente; una familia no central, por una multisección estable; una partícula compuesta, por una composición de registros genealógicos con frontera de enlace; una resonancia, por un polo del operador acoplado; un sector topológico, por su anillo de fusión y su representación de trenzado antes de toda carta de masa; una sección virtual, por su horizonte terminal de prolongación. Estas realizaciones son distintas y no deben ser aplanadas en una sola lista de números.

La magnitud externa nunca interviene en la selección de celda, ruta, firma, lector, coeficiente de torre ni autovector. La confrontación metrológica se efectúa después de construir y fijar la salida interna.

\Needspace{7\baselineskip}
\section{Espacio completo de estados de masa}
Sea \(\mathcal C_{81}\) el atlas celular. Cada celda admite cuatro orientaciones two--way, por lo que el espacio finito de rutas orientadas es

\[
 \mathcal R_{324}
 =\{(c,h_0,v_0):c\in\mathcal C_{81},\ (h_0,v_0)\in\{\pm1\}^2\}.
\]
La proyección celular, la relación de familia y la estratificación gruesa producen

\[
 |\mathcal C_{81}|=81,\qquad
 |\mathcal F_{56}|=56,\qquad
 |\mathcal M_{13}|=13,\qquad
 |\mathcal R_{324}|=324.
\]
El electrón ocupa el único punto fijo central. Las doce multisecciones no centrales son fibras finitas estables bajo la inversión celular y no admiten, en general, una sección puntual equivariante. Su objeto natural es, por tanto, un operador sobre la fibra y no un número elegido entre sus celdas. El sector nulo se incorpora mediante la sección cero de la recta de masa, no como valor de un carácter positivo.

El catálogo físico de salida distingue tres leptones cargados, tres realizaciones neutrínicas de interacción, seis sabores quark con sus órbitas de color, dos realizaciones de calibre nulas, dos bosones de calibre masivos, un sector escalar, dos familias compuestas, tres codominios topológicos y un sistema de secciones virtuales. Esta clasificación posterior no limita el número de estados compuestos o resonantes: composición y excitación generan familias adicionales dentro de sus codominios respectivos.

\Needspace{7\baselineskip}
\section{Trayectoria genealógica de una realización de masa}
La ley no comienza con una etiqueta de partícula ni con un número. Su orden es

\[
\boxed{
\begin{aligned}
\mathrm{APP}\to\mathrm{TRIT}\to\mathrm{TPK}
&\to\widetilde\Gamma^{\rm surv}
\to\Gamma^{\rm obs}
\to L_{108}
\to\sigma\in\mathbb Z^5\\
&\to\chi_0
\to\chi_{\rm full}
\to\widehat{\mathcal M}
\to\text{acoplamiento}\\
&\to\text{espectro/polo}
\to\mathcal L_M
\to\text{contraste}.
\end{aligned}}
\]
Aquí:

\begin{itemize}
\item APP dispone la incompatibilidad local, las hojas suma/producto y su orientación;
\item el TRIT conserva signo, umbral y apertura bidireccional;
\item el TPK evoluciona fase, memoria, acarreo, supervivencia y retorno;
\item una partícula es una sección persistente de una extensión, y su ruta observable es
la historia que deja en el registro genealógico;

\item la firma entera es aditiva;
\item el carácter convierte composición aditiva en una razón positiva;
\item las torres refinan el carácter mediante un único operador armónico global;
\item la recta de masa aporta el tipo dimensional;
\item medir es acoplarse: el acoplamiento determina qué subespacio y qué espectro son
físicamente accesibles, sin consultar la magnitud que después se medirá.

\end{itemize}
La masa no selecciona la ruta. La ruta persistente, el registro genealógico y el acoplamiento seleccionan el objeto espectral cuya masa se lee.

\Needspace{7\baselineskip}
\subsection{Componentes conservadas del estado enriquecido del TPK en la trayectoria de masa}
El estado que alimenta la ley de masas es

\[
 X_K=(w_6,w_{30},R_{36},g_K,B_K,p_K,c_K,\varepsilon_K,\omega_K),
\]
donde se conservan palabra corta, palabra prolongada, región, fase nonádica, memoria de bloque, prefijo, acarreo, hoja y devanado. La cadena

\[
 w_6\longrightarrow w_{30}\longrightarrow R_{36}
 \longrightarrow G_9
\]
no proyecta directamente una masa: construye la identidad persistente que después se lee.

\Needspace{7\baselineskip}
\subsection{Filtración nonádica con memoria}
Desde cada estado se obtiene una fibra de refinamiento formada por \(729\) subcilindros,

\[
 \{C_{K,j}:0\le j<729\}.
\]
La poda \(G_9\) satisface

\[
 G_9(C_{K,j})\in\{0,1\},\qquad
 \sum_{j=0}^{728}G_9(C_{K,j})=1.
\]
Si \(j_K\) es el superviviente,

\[
 P_{K+1}=729P_K+j_K.
\]
Tras nueve pasos,

\[
 g(K+9)=g(K),
\]
pero cambian, en general,

\[
 (B_{K+9},P_{K+9},\operatorname{pref}_{K+9},
 \operatorname{Witt}_{K+9})
 \ne
 (B_K,P_K,\operatorname{pref}_K,\operatorname{Witt}_K).
\]
Ésta es la forma ecuacional de la persistencia: retorna la fase, no el estado. La masa puede reconocer una partícula después del retorno porque el registro genealógico conserva la diferencia.

\Needspace{7\baselineskip}
\subsection{Transducción del acarreo nonádico al estado firmado dodecafásico}
Para una coordenada temporal \(\theta\), definimos

\[
 g_0(\theta)=\theta\bmod9,\qquad
 \beta(\theta)=
 \left\lfloor\frac{\widetilde\theta}{9}\right\rfloor\bmod12.
\]
Entonces

\[
 g_0(\theta+9)=g_0(\theta),\qquad
 \beta(\theta+9)=\beta(\theta)+1\pmod{12}.
\]
Así, el retorno nonádico conserva la fase y avanza un diente de memoria. Los 108 tiempos son

\[
 108=9\cdot12,
\]
de modo que cada registro genealógico contiene un ciclo completo de las dos escalas. Las 36.864 composiciones de acarreo describen las combinaciones locales de esta memoria bidireccional.

\Needspace{7\baselineskip}
\subsection{Completación y conservación de la unidad}
La frontera base mil se descompone como

\[
 999=729+270=729+243+27,
\]
mientras la completación es

\[
 1000=729+271=729+243+27+1,
\qquad
 271=270+1.
\]
El \(+1\) es el único vértice totalmente fronterizo y realiza el acarreo que devuelve la unidad. La partición \(729/271\) distingue núcleo activo y borde mediador; no es una aproximación áurea. La misma conservación permite que una ruta cambie de bloque sin perder su identidad normalizada.

\Needspace{7\baselineskip}
\subsection{Teorema del atlas orientado completo}
Cada celda \(c\in\mathcal C_{81}\) posee cuatro orientaciones

\[
 (h_0,v_0)\in\{\pm1\}^2.
\]
El mapa

\[
 (c,h_0,v_0)\longmapsto\Gamma(c,h_0,v_0)
\]
es biyectivo sobre \(\mathcal R_{324}\). En consecuencia,

\[
 |\mathcal R_{324}|=4|\mathcal C_{81}|=324.
\]
Cada ruta posee 108 entradas de registro genealógico, de modo que el censo temporal total es

\[
 81\cdot108=8748.
\]
La cuádrupla

\[
 (\text{fila semilla},\text{columna semilla},h_0,v_0)
\]
determina de manera única la ruta, su palabra \(w_6\), su palabra dodecafásica, sus retornos, memoria, acarreo, firma histórica y lectores. Por tanto, las 324 rutas no sólo pueden enumerarse: pueden reconstruirse y unirse exactamente, campo por campo, sin utilizar una masa.

\Needspace{7\baselineskip}
\subsection{Cocientes de familia y multisección}
El cociente de las 81 celdas por historia de ruta produce 56 familias. Sus multiplicidades se distribuyen como

\[
 41\times1,\quad10\times2,\quad2\times3,\quad1\times4,\quad2\times5,
\]
y

\[
 41+20+6+4+10=81.
\]
La estratificación en trece multisecciones posee rangos

\[
 1,\quad2,\quad4,\quad6,\quad8,\quad12,\quad16
\]
con multiplicidades

\[
 1,\quad2,\quad3,\quad2,\quad3,\quad1,\quad1,
\]
y suma de rangos

\[
 1+4+12+12+24+12+16=81.
\]
El electrón ocupa el bloque de rango uno. Los demás sectores viven en bloques de rango mayor y exigen un proyector de representación, no una celda nominal.

\Needspace{7\baselineskip}
\subsection{Firma, lectores y memoria por ruta}
Para cada \(\Gamma\in\mathcal R_{324}\) se calculan:

\[
 \sigma_\Gamma\in\mathbb Z^5,\qquad
 K_D(\Gamma),\qquad K_\Omega(\Gamma),
\]
junto con primeros retornos de clase, órbita, celda y estado cinemático; separación de prefijos; deuda de acarreo; ambigüedad; tríadas estabilizadas; hoja y devanado. La familia \(K_D\) presenta 14 perfiles; \(K_\Omega\), nueve; juntos forman 40 pares y 18 coincidencias. Esta diversidad es la información fina que permite pasar de una fibra gruesa a un operador físico.

\Needspace{7\baselineskip}
\section{Del defecto APP a la firma entera}
\Needspace{7\baselineskip}
\subsection{Defecto primordial y despliegue}
En el centro APP,

\[
 B_c=9P_++5P_-,\qquad 9-5=4.
\]
El salto local primordial es 4. El defecto global \(54\) pertenece a un despliegue posterior \(2\to6\to54\); no debe retroproyectarse sobre el centro como si fuese la misma operación.

\Needspace{7\baselineskip}
\subsection{\(12\) eventos y divisor dodecafásico}
La coordenada orientada no recibe arbitrariamente un factor \(1/6\). La cadena es

\[
12\ \text{eventos}
\to s_C^{\rm tot}
\to 1+5+6
\to s_C^{(12)}=\frac{s_C^{\rm tot}}6
\to W_\pm=6n_A\pm s_C
\to\operatorname{SNF}(1,12)
\to\chi_{\rm mass}.
\]
El \(1/6\) es el cociente orientado del cierre de doce eventos. No es el residuo (-1/12), no es la pseudoinversa de una valencia y no es la incidencia hexada--octada.

\Needspace{7\baselineskip}
\subsection{Construcción de la firma entera completa desde el registro de ruta}
El registro genealógico CT--108 produce

\[
 \sigma(\Gamma)
 =\bigl(n_A,s_C,k_{\Delta_4},b_{120},b_\zeta\bigr)\in\mathbb Z^5.
\]
Las coordenadas tienen procedencias distintas:

\begin{itemize}
\item \(n_A\): duración/persistencia dominante;
\item \(s_C\): orientación neta de la bicapa, dividida por seis tras el cierre
dodecafásico;

\item \(k_{\Delta_4}\): canal de torsión mínima;
\item \(b_{120}\): coeficiente de evento del registro;
\item \(b_\zeta\): autocorrelación de paso tres.
\end{itemize}
Se define

\[
 \zeta_{270}:=135\Delta_4^2,
\]
y se mantiene separado de

\[
 s_{270}=q_-^{270}-q_+^{270}.
\]
Las 324 rutas orientadas se unen biyectivamente con sus cinco coordenadas históricas mediante la identidad canónica de ruta. Esta unión usa únicamente estructura genealógica y no consulta ninguna magnitud metrológica.

\Needspace{7\baselineskip}
\subsection{Ángulo directo, contraángulo y canales conjugados}
La coordenada angular directa procede de la constante de estructura fina ya construida:

\[
 A=1000\alpha_{\rm HMT}
 =7.297352569283800997285105472380663\ldots{}^\circ.
\]
El contraángulo no se obtiene de una masa ni del funcional de Barbero--Immirzi. Sea

\[
 D_A=\exp\!\left(-\frac{100\pi}{9}\alpha_{\rm HMT}\right),
\]
\[
 \mathcal C_\pi(\alpha_{\rm HMT})
 =169\alpha_{\rm HMT}^{6}
 +\frac{D_A\alpha_{\rm HMT}^{7}}
        {1-D_A\alpha_{\rm HMT}},
\]
\[
 y_*=\frac{\pi}{90}(H_5+\alpha_{\rm HMT})
      -\mathcal C_\pi(\alpha_{\rm HMT}),
 \qquad
 C^*=\frac{180}{\pi}y_*.
\]
Por tanto,

\[
 C^*=2.123738338968645724261951380393\ldots{}^\circ.
\]
La coordenada \(A\) es el modo angular común y \(C^*\), el modo diferencial orientado. Su paso a la carta arquimediana es

\[
 x=A_{\rm rad}=\frac{\pi A}{180},
 \qquad
 y=C^*_{\rm rad}=\frac{\pi C^*}{180}.
\]
En el álgebra real generada por la involución \(R=\operatorname{diag}(1,-1)\), con \(P_\pm=(I\pm R)/2\), el bloque biangular es

\[
 B_{A,C^*}=AI+C^*R
 =(A+C^*)P_++(A-C^*)P_-.
\]
Ángulo más contraángulo y ángulo menos contraángulo son, así, los dos autovalores de un único objeto. La exponenciación produce los canales

\[
 q_+=e^{-x-y},\qquad q_-=e^{-x+y}.
\]
La carta es invertible:

\[
 x=-\frac12\log(q_+q_-),
 \qquad
 y=\frac12\log\frac{q_-}{q_+}.
\]
La transformación \(C^*\mapsto-C^*\) conserva los proyectores e intercambia los canales. La coordenada de retorno satisface, después de la construcción regional,

\[
 \eta_{\rm ret}=\frac{C^*}{2}-\alpha_{\rm HMT}\,{\rm deg},
\]
sin constituir una segunda regla de selección de \(C^*\).

\Needspace{7\baselineskip}
\subsection{Carácter central}
Con

\[
 \Theta_0=
 \left(A_{\rm rad},\frac{C^*_{\rm rad}}6,
 \Delta_4,s_{120},\zeta_{270}\right),
\]
el carácter central es

\[
 \chi_0(\sigma)=
 \exp\!\left(
 n_AA_{\rm rad}+\frac{s_C}{6}C^*_{\rm rad}
 +k_{\Delta_4}\Delta_4+b_{120}s_{120}+b_\zeta\zeta_{270}
 \right).
\]
Por aditividad de la firma,

\[
 \chi_0(\sigma_1+\sigma_2)=\chi_0(\sigma_1)\chi_0(\sigma_2).
\]
La unidad física no entra en esta exponencial: \(\chi_0\) es un automorfismo positivo adimensional de la recta de masa.

\Needspace{7\baselineskip}
\section{Jerarquía armónica global de torres}
Con las coordenadas (x=\allowbreak{}A\_\{\textbackslash{}rm rad\}) e (y=\allowbreak{}C\textasciicircum{}*\_\{\textbackslash{}rm rad\}) ya construidas, sean

\[
 q_\pm=\exp\!\left[-\frac{\pi}{180}(A\pm C^*)\right]
 =e^{-x\mp y}.
\]
En una convención global coherente,

\[
 T=e^{-xI-yR},\qquad
 c_n=\operatorname{Tr}(T^n)=q_-^n+q_+^n,
\]
\[
 s_n=-\operatorname{Tr}(RT^n)=q_-^n-q_+^n.
\]
Equivalente y operatorialmente,

\[
 c_n=2e^{-nx}\cosh(ny),\qquad
 s_n=2e^{-nx}\sinh(ny).
\]
La involución \(C^*\mapsto-C^*\) deja fijos los proyectores y permuta los coeficientes espectrales, o, equivalentemente, los canales asociados. De aquí se obtienen

\[
 c_{m+n}=\frac{c_mc_n+s_ms_n}{2},\qquad
 s_{m+n}=\frac{s_mc_n+c_ms_n}{2},
\]
\[
 c_n^2-s_n^2=4e^{-2nx},
\]
y la recurrencia común de orden dos. Las alturas admisibles forman

\[
 \mathfrak S=\langle90,120\rangle
 =\{90,120\}\cup\{30r:r\ge6\}.
\]
Los valores \(90,120,180,210,240,270,360,420\) son testigos genealógicos, no la torre completa. Antes del cociente, la altura 360 conserva dos historias:

\[
 k[\mathfrak S_0]
 \simeq k[X_{90},Y_{120}]/(X_{90}^4-Y_{120}^3).
\]
La elevación general se escribe

\[
 \chi_{\rm full}(\sigma,\eta)
 =\chi_0(\sigma)
  \exp\!\left(\sum_{h\in\mathfrak S}\eta_hs_h\right),
\]
con

\[
 N_{120}=b_{120}+\eta_{120}.
\]
La suma numérica no borra la doble procedencia: \(b_{120}\) viene del registro y \(\eta_{120}\) del refinamiento espectral. Tampoco puede absorber \(b_\zeta\zeta_{270}\) en \(\eta_{270}s_{270}\).

\Needspace{7\baselineskip}
\section{Rama absoluta: acción, reloj y recta de masa}
El lector determinantal produce

\[
 \operatorname{SNF}(H_4)=\operatorname{diag}(1,2,2,4),\qquad
 |\operatorname{coker}H_4|=16,
\]
\[
 \Sigma_H=\varphi\alpha_{\rm HMT}^{16},\qquad
 \operatorname{ord}_{10}(\Sigma_H)=-34.
\]
Las dos secciones orientadas de una sola recta de acción son

\[
 \hbar_{\rm pre}=H_5\,10^{-34}\mathcal U_S,
 \qquad
 \hbar_{\rm ret}=\eta_{\rm ret}\,10^{-34}\mathcal U_S.
\]
Su carta aditiva fina y su carta multiplicativa son

\[
 \delta_{2w}=H_5-\eta_{\rm ret},\qquad
 R_{\rm act}=\frac{H_5}{\eta_{\rm ret}}>0.
\]
No son dos constantes de Planck. La década común fija el orden absoluto; el cociente transporta información relativa de ida/retorno.

Declarado \(t_0\in\mathcal L_T^+\), CT--108 fija

\[
 \omega_0=\frac{2\pi}{108t_0},\qquad
 \varepsilon_{\rm clk}=\hbar_{\rm ret}\omega_0,
 \qquad
 m_{\rm clk}=\varepsilon_{\rm clk}c_{\rm int}^{-2}.
\]
El carácter actúa por homotecia sobre \(\mathcal L_M\); no se convierte en \(\hbar\). En cocientes de masas, todo factor común se cancela:

\[
 \frac{m_Y}{m_X}
 =\chi_{\rm full}(\sigma_Y-\sigma_X,\eta_Y-\eta_X)
 R_{\rm act}^{\kappa_Y-\kappa_X}.
\]
Por ello \(10^{-34}\) fija la escala absoluta, pero no corrige un residual relativo. Una modificación relativa sólo puede proceder de la ruta, el reloj específico, el lector o las torres.

Queda una distinción necesaria. HMT construye una sección interna una vez declarado \(t_0\); la coordenada MeV exige una carta. Para cerrar la carta absoluta sin una convención externa debe construirse

\[
 \Gamma\longmapsto t_\Gamma
 \quad\text{o}\quad
 \Gamma\longmapsto\omega_\Gamma.
\]
La sección 10 define un lector de reloj basado en retornos y holonomía. Su normalización absoluta queda determinada cuando el operador de ruta fija una sección positiva de la recta temporal.

\Needspace{7\baselineskip}
